\section{Corpus Construction and Validation}
\label{sec:method}

\subsection{Game Selection}

The 150-game corpus covers BGG weight 1.07--4.70, publication years 1876--2024,
and player counts 1--8. The corpus is stratified across mechanism types:
party/social games (weight $< 2.0$, $n=42$), medium-weight Euros (2.0--3.0, $n=48$),
heavy Euros ($> 3.0$, $n=35$), abstract strategy ($n=8$), originals designed within
the TIGRIS factory ($n=17$).

\subsection{Axis Definitions and Scoring Rubric}

The 32 axes are organized in four bands (A--D) plus four new axes (A29--A32)
introduced in this work.
Each axis is defined with \emph{anchor descriptors} at scores 0, 5, and 10,
specifying games that exemplify each anchor point.
For example, A3 (Interaction) anchors: 0 = ``Wingspan base game (parallel tableau optimization)'';
5 = ``Catan (indirect competition via resource blocking)'';
10 = ``Crokinole (the game is the physical contest)''.

The four new axes are:
\begin{itemize}
  \item \textbf{A29 Forgiveness Curve} (0 = punishing snowball $\to$ 10 = forgiving recovery)
  \item \textbf{A30 Strategy Ergodicity} (0 = single optimal path/solved $\to$ 10 = many equally-valid strategies)
  \item \textbf{A31 Failure Texture} (0 = boring loss, just a point deficit $\to$ 10 = dramatic narrative failure)
  \item \textbf{A32 Information Architecture Type} (0 = fully transparent $\to$ 10 = inverted visibility)
\end{itemize}

\subsection{Two Scoring Protocols}

\textbf{Full-pipeline (Parliament procedure), 69 games:}
An adversarial multi-persona review in which 8 evaluator personas draft stakes on axes,
argue for/against each axis via narrated playthrough argument records, and classify each
stake as earned/refuted/contested. This produces Parliament-confirmed scores grounded in
explicit argument records.

\textbf{Fast-track independent assessment, 81 games:}
A single evaluator assigns honest 0--10 scores to all 32 axes independently,
without adversarial procedure. Scores are based on the evaluator's direct knowledge
of the game's rules and mechanisms, cross-referenced against published rules and BGG data.
The fast-track protocol is designed to avoid the \emph{Parliament drafting bias}
identified in preliminary analysis: the adversarial procedure inflates scores because
evaluators draft their preferred axes and then find evidence for them,
producing a general quality factor that explains 81\% of raw-score variance.
Fast-track scores reduce this to $\approx$73\% and within-game normalization to 50\%.

\subsection{Inter-Rater Reliability Study Design}

We validate the instrument via a held-out reliability study.
The 20-game held-out sample was selected to represent the full design space:
5 light games (weight $<$2.0), 5 medium Euros (2.0--3.0), 5 heavy games ($>$3.0),
3 originals, and 2 abstract strategy games.

Five rater profiles were recruited to represent distinct expert communities:
\begin{enumerate}
  \item \textbf{Marcus} — Euro game designer, 15 years experience, Feld/Rosenberg tradition;
  \item \textbf{Priya} — Cognitive scientist, studies reasoning and decision-making in games;
  \item \textbf{Jin} — Accessibility researcher, party game enthusiast;
  \item \textbf{Elena} — BGG power user, has rated $>$800 games;
  \item \textbf{Theo} — Competitive abstract strategy player (Chess, Go, Hive).
\end{enumerate}

Raters scored all 20 games on all 32 axes independently, with no access to each other's
scores. Each rater received the axis definitions and anchor descriptors but not
the game-specific scores from the TIGRIS corpus.
We compute Cronbach's $\alpha$~\citep{cronbach1951coefficient} treating each rater as an
item and each game as an observation, per axis.
